\section{结论与展望}
\label{sec:conclusion}
%------------------------------------------------------------------------------

自 Feynman 五十多年前提出部分子模型以来，我们对质子部分子结构的理解已大为推进。一方面，SLAC、DESY、CERN、Fermi 实验室、JLab、BNL 等世界各地装置上开展的大量高能实验，使我们得以在不同能量与极化下探测强子结构的方方面面。另一方面，许多部分子观测量被并行提出，为质子结构提供多维描述，包括共线 PDF、TMDPDF、GPD、部分子 DA、LFWF 等。

虽然带 RG 改进的 QCD 因子化定理允许我们经由这些部分子观测量与实验观测量的联系来提取它们，但非常期望能通过格点 QCD 这类 {\it ab initio} 计算从理论上预言它们。由于模拟实时动力学的困难，这一方向的进展相当缓慢。然而，自几年前 LaMET 被提出以来，情况已经改变——它提供了从第一性原理计算部分子物理的系统可改进方法。

本文概述了 LaMET 的形式体系及其在格点 QCD 与其他非微扰方法可及观测量上的应用。通过考察束缚态强子结构的参考系依赖，我们解释了 IMF 物理（或部分子物理）如何作为质子结构的 EFT 描述自然产生。这种 EFT 描述在 SCET 与 LFQ 中表述得最自然，但质子矩阵元的实际非微扰计算一直困难。LaMET 实际上提供了实现 LFQ 所需的东西：在大动量态中构造恰当的准部分子观测量，并通过因子化把它们匹配到 LF 上的真正部分子观测量。就 PDF 而言，前者对应有限动量分布，其跑动由动量 RGE 控制；后者对应 IMF PDF，其跑动由通常的 RGE 控制。应当指出，LaMET 是一个非常一般的框架，适用于任何非微扰方法（欧氏的虚时间方法或 Minkowski 的实时间方法）计算的大动量物理量。而且，给定一个大动量态，同一部分子物理可以由构成普适类的不同准可观测量确定。

接着我们展示了如何在实践中计算部分子观测量，重点放在共线 PDF、GPD、DA、TMDPDF 与 LFWF 上。我们还讨论了质子自旋结构，说明如何用同样的途径获得部分子对质子自旋的贡献。最后我们总结了迄今用 LaMET 开展的格点研究：一方面证明 LaMET 是计算质子部分子结构的有希望的途径；另一方面也清楚表明，要使格点结果对唯象学产生可观影响，仍需大量改进。

我们以对未来格点计算改进的几点评论作结。
%
关于其中某些问题（例如连续、无穷体积与物理 $\pi$ 介子质量极限）更系统的讨论，推荐参阅~\cite{Alexandrou:2019lfo}。

\begin{itemize}

	\item 大强子动量。LaMET 的未来在于更大的动量，这天然要求更小的格距。因此，在百亿亿次计算中应对大动量与小格距带来的挑战——如激发态污染或拓扑荷冻结问题——将至关重要。

	\item 重正化。如 \sec{renorm-npr} 所讨论，Wilson 线算符的质量重正化是可取的：它规范不变，且不引入额外的高扭度效应或长距离处的大统计误差。然而其到 $\MS$ 方案的匹配、尤其是重正子模糊性，仍需解决才能完整系统地应用。此外，非常期待包含上述优点的替代方案。

	\item 高阶微扰匹配。在当前的 LaMET 计算中，一圈微扰匹配带来了可观的修正。要控制该流程的系统误差，需要更高阶的匹配核。

	\item 幂次修正。当强子动量不太大或 $x$ 接近 0 或 1 时幂次修正很重要。迄今为止朝模型无关地确定幂次修正进展甚微。一种权宜策略是在实施匹配与靶质量修正之后外推到 $P^z\to\infty$ 极限；但根本出路在于 \sec{others-ht} 讨论过的高扭度分布的格点计算。

\end{itemize}

上述关于系统误差的讨论是普适的，适用于所有准可观测量。过去几年丰富的理论发展为用 LaMET 计算广泛的部分子观测量铺平了道路。随着计算资源的迅速增长以及新技术与新算法的发展，我们预期上述系统误差未来会逐步得到控制。这对确立 LaMET 作为计算部分子物理的系统方法、并让格点计算在 EIC 时代扮演关键角色十分重要。


%------------------------------------------------------------------------------
